\chapter{Conclusiones}
\label{chap:conclusiones}

\section{Conclusión}

El objetivo general era construir y validar un banco de ensayos para un mecanismo de elevación con
caracola de radio variable, y usarlo para contrastar el modelo del mecanismo. El banco está
construido y funciona; el contraste ha dado un resultado distinto del esperado, y de ahí salen las
conclusiones más interesantes del trabajo.

\subsection{Sobre el banco de ensayos}

\begin{enumerate}
  \item El sistema cumple los requisitos de temporización: muestrea a \SI{20.0}{\hertz} con una
        desviación típica del periodo de \SI{0.06}{\milli\second}, y el trabajo del lazo ocupa el
        \SI{22}{\percent} del periodo. Llegar a ese determinismo exigió abandonar la versión de doble
        núcleo en favor de un lazo de periodo fijo con un único dueño de cada bus.
  \item El ciclo de ensayo B$\rightarrow$A$\rightarrow$B, con los extremos reconocidos por los
        finales de carrera y no por la posición integrada, repite el recorrido dentro del
        \SI{0.6}{\percent} y cierra sobre el punto de partida. Esa decisión hace que el banco sea
        utilizable pese a la incertidumbre de la medida de posición.
  \item Guardar los ensayos en cuentas brutas, con las constantes en la cabecera, ha permitido
        reanalizar por completo dos campañas después de cambiar el modelo y la calibración, sin
        repetir un solo ensayo. Es la decisión de diseño que más ha rendido.
  \item Unificar el modelo en una sola implementación compartida por el simulador y el visor
        eliminó una fuente real de error: las páginas habían divergido hasta calcular con
        geometrías distintas.
\end{enumerate}

\subsection{Sobre el modelo del mecanismo}

\begin{enumerate}\setcounter{enumi}{4}
  \item La posición de la salida del cable respecto a la articulación no es un detalle de dibujo.
        Con la salida en la vertical, la tensión resulta proporcional a la longitud libre de cable,
        un resultado simple y atractivo; con la salida desplazada, esa cancelación desaparece y la
        tensión pasa a depender de dos ángulos. El modelo del banco trabaja ya con el caso general.
  \item El perfil de par constante es el que minimiza el par de pico para un recorrido y un giro de
        tambor dados. La demostración es independiente de la geometría: el trabajo está fijado, luego
        el máximo del par no puede ser menor que su media.
\end{enumerate}

\subsection{Sobre la validación experimental}

\begin{enumerate}\setcounter{enumi}{6}
  \item \textbf{Medir sólo la fuerza no permite validar el modelo.} Con la célula sin calibrar, el
        ajuste incluye la sensibilidad como incógnita y el coeficiente de determinación es el
        cuadrado de la correlación: no distingue el signo. Tres hipótesis de tambor muy distintas
        ---cilindro, caracola y caracola invertida--- ajustan la señal de la célula con
        $R^2$ entre \num{0.93} y \num{0.97}. Un buen ajuste no es una validación.
  \item \textbf{El par del accionamiento sí decide}, porque no depende de ninguna calibración
        pendiente. Y al usarlo, la comparación es concluyente: un tambor cilíndrico obligaría al par
        a bajar un \SI{22}{\percent} a lo largo del tramo con carga, mientras que el par medido sube
        un \SI{56}{\percent}. \textbf{El tambor no se comporta como un cilindro}: su radio efectivo
        crece un \SI{93}{\percent} a lo largo del recorrido.
  \item \textbf{La caracola, sin embargo, no compensa: sobrecompensa.} Para mantener el par constante
        bastaría con que el radio creciera un \SI{24}{\percent}. Con el crecimiento real, el par
        queda peor repartido que con un cilindro equivalente, y además en el sentido contrario al que
        describe la configuración del banco.
  \item Las pérdidas del mecanismo son de tipo Coulomb: constantes a lo largo del recorrido y del
        orden del \SI{19}{\percent} del par gravitatorio.
  \item El resultado anterior no cierra la validación, y es importante decirlo con claridad: mientras
        no se midan las cotas del mecanismo y no se calibre la célula con masas patrón, cualquier
        juego de parámetros es una descripción de los datos y no una verificación del modelo. El
        trabajo deja el procedimiento montado y el diagnóstico hecho.
\end{enumerate}

\subsection{Grado de cumplimiento de los objetivos}

\begin{longtable}{@{}L{0.07\textwidth}L{0.58\textwidth}L{0.25\textwidth}@{}}
\toprule
\textbf{Obj.} & \textbf{Resultado} & \textbf{Estado} \\
\midrule
\endhead
O1 & Modelo estático general, perfil ideal de par constante y demostración de que minimiza el par de pico (capítulo~\ref{chap:modelo}). & Cumplido \\
O2 & Placa adaptadora con RS-485, entradas y salidas y dos zócalos mikroBUS; revisión A analizada y corregida en la revisión B (capítulo~\ref{chap:hardware}). & Cumplido; falta versionar la revisión B \\
O3 & Firmware determinista con Modbus, célula, \emph{homing}, ciclo B$\rightarrow$A$\rightarrow$B y protecciones (capítulo~\ref{chap:firmware}). & Cumplido \\
O4 & Protocolo binario con CRC-8, telemetría por lotes y compatibilidad entre versiones (capítulo~\ref{chap:comunicaciones}). & Cumplido \\
O5 & Servidor y tres páginas web: monitor, configuración con simulador y visor de ensayos (capítulo~\ref{chap:software}). & Cumplido \\
O6 & Caracterización del banco y contraste cilindro/caracola sobre dos campañas (capítulo~\ref{chap:resultados}). & Cumplido; validación abierta a falta de cotas y calibración \\
O7 & Problemas y soluciones, planificación, presupuesto e impacto (capítulos~\ref{chap:problemas} y~\ref{chap:gestion}). & Cumplido \\
\bottomrule
\end{longtable}

\section{Desarrollos futuros}

\subsection{Para cerrar la validación}
\begin{itemize}
  \item Medir las cotas del mecanismo y el perfil real de la caracola (radios, sentido y barrido), y
        sustituir con ellas los parámetros ajustados.
  \item Calibrar la célula de carga con masas patrón y estimar su incertidumbre; con ello, los
        valores de fuerza pasarían de orientativos a trazables.
  \item Confirmar la escala del registro de par del accionamiento para poder expresar los pares en
        \si{\newton\metre}.
  \item Eliminar o medir la holgura inicial del cable, para que el recorrido con carga coincida con
        el recorrido completo.
  \item Repetir la campaña y rehacer el contraste cilindro frente a caracola, que el banco ya
        automatiza.
\end{itemize}

\subsection{Sobre el mecanismo}
\begin{itemize}
  \item Fabricar, por ejemplo por impresión 3D, un tambor cilíndrico y otro con el perfil ideal
        calculado en la sección~\ref{chap:modelo}, y ensayar los tres en el mismo banco. Sería la
        comparación definitiva, con el mecanismo como única variable.
  \item Revisar el sentido de montaje de la caracola actual, que según la medida trabaja al contrario
        de lo que describe la configuración.
  \item Ensayar con varias masas para comprobar la linealidad del modelo con la carga.
\end{itemize}

\subsection{Sobre el sistema de medida}
\begin{itemize}
  \item Leer la posición del codificador del accionamiento por Modbus en lugar de integrar la
        velocidad.
  \item Leer velocidad y par en una sola transacción Modbus (registros 16385 a 16387), lo que
        ahorraría unos \SI{3}{\milli\second} por periodo y permitiría subir la frecuencia de muestreo.
  \item Enviar en cada muestra la posición del firmware y la lectura sin filtrar de la célula, y
        sincronizar la lectura con la señal de interrupción del acondicionador.
  \item Corregir la codificación de la ganancia y los bits de estado del ZSC31014
        (sección~\ref{sec:prob-estado}).
\end{itemize}

\subsection{Sobre la seguridad y la explotación}
\begin{itemize}
  \item Cablear los finales de carrera y una seta de emergencia a la entrada STO del accionamiento,
        de modo que la parada no dependa del software.
  \item Versionar y fabricar la revisión B de la placa, con los ajustes de resistencias propuestos.
  \item Hacer que el servidor escuche sólo en la máquina local o exija autenticación, y añadir
        pruebas automáticas del protocolo.
\end{itemize}

\section{Valoración personal}

\completar{Valoración personal: qué se ha aprendido, qué dificultades han sido más formativas y qué
se haría de otra manera.}
